% =====================================================================
%  Chuong01_TongQuan.tex — Chương 1: Tổng quan về đề tài
%
%  Đánh số chương bằng \chapter (không dùng \chapter*) để đồng bộ
%  với các chương còn lại và khớp với mô tả trong Mục "Bố cục báo cáo".
% =====================================================================

\chapter{Tổng quan về đề tài}
\label{introduction}

\section{Bối cảnh và động lực nghiên cứu}

Sự phát triển của môi trường thông tin trực tuyến đã làm thay đổi đáng kể cách người dùng Việt Nam thu nhận và truyền tải thông tin. Các nền tảng như Facebook, TikTok, Zalo cùng với những trang tin điện tử cho phép nội dung được phổ biến đến số lượng lớn người dùng trong thời gian ngắn. Tuy nhiên, tốc độ lan truyền cao cũng khiến những nội dung chưa được kiểm chứng dễ dàng tiếp cận cộng đồng, từ đó làm gia tăng nguy cơ thông tin sai lệch được tiếp nhận và chia sẻ rộng rãi.

Khối lượng nội dung ngày càng lớn khiến việc xác minh thủ công từng bài viết trở nên khó duy trì ở quy mô thực tế. Nếu quá trình kiểm tra phụ thuộc hoàn toàn vào con người, chi phí về thời gian và nhân lực sẽ tăng đáng kể khi số lượng bài viết được tạo ra mỗi ngày ngày càng nhiều. Đây là động lực để nghiên cứu các phương pháp tự động hóa quá trình phát hiện thông tin sai lệch, đặc biệt thông qua việc kết hợp các kỹ thuật học máy và học sâu trong xử lý văn bản tiếng Việt.

\section{Mục tiêu của đồ án}

Từ nhận định trên, đồ án đặt ra bốn mục tiêu cụ thể:

\begin{itemize}[leftmargin=2em]
  \item Xây dựng và đánh giá song song bốn mô hình phân lớp cho văn bản tiếng Việt, gồm: Hồi quy Logistic (LR), Máy vectơ hỗ trợ (SVM), mạng LSTM hai chiều (BiLSTM) và mô hình ngôn ngữ tiền huấn luyện PhoBERT.
  \item Khảo sát tác động của bước phân đoạn từ --- một đặc thù riêng của tiếng Việt --- đến chất lượng dự đoán của mỗi mô hình.
  \item Đánh giá đa chiều: kiểm định thống kê, kiểm tra chéo và phân tích khả năng giải thích của mô hình.
  \item Từ các kết quả thu được, chỉ ra phương pháp phù hợp nhất cho bài toán phát hiện thông tin sai lệch trên tiếng Việt.
\end{itemize}
\section{Phạm vi nghiên cứu}

Nghiên cứu tập trung vào việc nhận diện tin giả trong văn bản tiếng Việt dưới dạng bài toán phân loại hai lớp, tương ứng với \textit{thật} và \textit{giả}. Dữ liệu đầu vào được hình thành từ hai thành pcủa bài viết là tiêu đề và nội dung, qua đó cung cấp thông tin cho quá trình phân loại. Tập dữ liệu được thu thập từ các bài đăng công khai trên một số nền tảng mạng xã hội phổ biến tại Việt Nam và từ các bài viết xuất hiện trên các trang báo điện tử.

Các dạng dữ liệu và nhiệm vụ nằm ngoài phạm vi nghiên cứu gồm nội dung hình ảnh, video và các tín hiệu lan truyền trên môi trường mạng. Đồ án cũng không xem xét bài toán phát hiện tin giả trong nhiều ngôn ngữ, đồng thời không triển khai các hình thức phân loại đa lớp hoặc đa nhãn.
\section{Cấu trúc báo cáo}

Nội dung của đồ án được tổ chức theo trình tự từ việc xác định bài toán đến xây dựng phương pháp, thực nghiệm và phân tích kết quả. Báo cáo gồm năm chương và phần kết luận, với nội dung chính của từng phần được trình bày như sau:

\begin{itemize}[leftmargin=2em]

    \item \textbf{Chương 1}: Đặt nền tảng cho đề tài thông qua việc trình bày vấn đề nghiên cứu, mục tiêu hướng đến và phạm vi triển khai. Nội dung chương cũng cung cấp cái nhìn tổng quan về bố cục của các phần tiếp theo trong báo cáo.

    \item \textbf{Chương 2}: Cung cấp những kiến thức nền tảng cho nghiên cứu, bao gồm xử lý ngôn ngữ tự nhiên tiếng Việt, những phương pháp học máy và kiến trúc học sâu được sử dụng. Các khái niệm liên quan đến biểu diễn văn bản và mô hình ngôn ngữ cũng được nhắc tới trong chương.

    \item \textbf{Chương 3}: Trình bày quá trình chuẩn bị dữ liệu cho thực nghiệm, từ mô tả nguồn dữ liệu, các bước tiền xử lý đến phương pháp phân chia dữ liệu thành các tập Train, Validation và Test.

    \item \textbf{Chương 4}: Trình bày quá trình triển khai bốn mô hình phát hiện tin giả, từ việc xây dựng đầu vào và thiết lập quá trình huấn luyện đến các phương pháp xử lý sự chênh lệch giữa các lớp dữ liệu.

    \item \textbf{Chương 5}: Phân tích các kết quả thu được qua thực nghiệm và đánh giá hiệu năng của từng mô hình. Chương cũng xem xét các trường hợp dự đoán sai, thực hiện kiểm định thống kê và tiến hành các thí nghiệm ablation để đánh giá tác động của từng thành phần trong hệ thống.

    \item \textbf{Phần kết luận}: khái quát những kết quả chính của nghiên cứu, đồng thời thảo luận các hạn chế còn tồn tại và đề xuất những hướng phát triển có thể thực hiện trong tương lai.

\end{itemize}
